% The instructor's eight-page handwritten notes end with the above limit.
% This short section gives the concrete free spinor solutions used in the derivation.
\section{Free Dirac Spinors}
For a plane wave $\psi=u(\vect p)e^{\ii(\vect p\cdot\vect r-Et)/\hbar}$ with $E>0$, the upper and lower components obey
\[
(E-mc^2)\varphi=c(\vect\sigma\cdot\vect p)\chi,\qquad
(E+mc^2)\chi=c(\vect\sigma\cdot\vect p)\varphi.
\]
Consequently $\chi=c(\vect\sigma\cdot\vect p)\varphi/(E+mc^2)$. For a normalized two-component spin basis $\xi_s$ with $s=\uparrow,\downarrow$, choose
\[
\boxed{u_s(\vect p)=\sqrt{\frac{E+mc^2}{2E}}
 \begin{pmatrix}\xi_s\\\displaystyle\frac{c\vect\sigma\cdot\vect p}{E+mc^2}\xi_s\end{pmatrix},\qquad
 u_s^\dagger u_{s'}=\delta_{ss'}.}
\]
This particular normalization is to unit probability density, $u_s^\dagger u_s=1$; alternative Lorentz-covariant normalizations are common in quantum field theory. The lower two components are suppressed as $|\vect p|/(2mc)$ in the nonrelativistic limit. Negative-energy solutions are obtained from the $E<0$ branch of the same $4\times4$ Hamiltonian and are interpreted consistently in a quantized Dirac field.